% !TEX program = xelatex
\documentclass[10pt, aspectratio=169]{beamer}
\usetheme{metropolis}

\usepackage{fontspec}
\setsansfont{Noto Sans CJK KR}[
  Path=/usr/share/fonts/opentype/noto/,
  Extension=.ttc,
  UprightFont=NotoSansCJK-Regular,
  BoldFont=NotoSansCJK-Bold,
  FontIndex=1, BoldFontIndex=1
]
\setmonofont{Noto Sans Mono CJK KR}[
  Path=/usr/share/fonts/opentype/noto/,
  Extension=.ttc,
  UprightFont=NotoSansCJK-Regular,
  BoldFont=NotoSansCJK-Bold,
  FontIndex=6, BoldFontIndex=6,
  Scale=0.8
]
\usepackage{graphicx}
\usepackage{booktabs}
\usepackage{tabularx}
\usepackage{array}
\usepackage{xcolor}

\definecolor{mainblue}{HTML}{1976D2}
\definecolor{darkblue}{HTML}{0D47A1}
\definecolor{lightblue}{HTML}{E3F2FD}
\definecolor{tipgreen}{HTML}{E8F5E9}
\definecolor{warnamber}{HTML}{FFF8E1}
\definecolor{keymsg}{HTML}{FCE4EC}
\definecolor{cogreen}{HTML}{2E7D32}
\definecolor{cored}{HTML}{C62828}

\setbeamercolor{palette primary}{bg=mainblue}
\setbeamercolor{frametitle}{bg=mainblue, fg=white}
\setbeamercolor{title}{fg=darkblue}
\setbeamercolor{alerted text}{fg=cored}

\graphicspath{{images/}}
\newcolumntype{L}[1]{>{\raggedright\arraybackslash}p{#1}}
\newcolumntype{C}[1]{>{\centering\arraybackslash}p{#1}}
\newcommand{\COtwo}{CO\textsubscript{2}}
\newcommand{\tld}{\raise.17ex\hbox{$\scriptstyle\sim$}}

% ── 한글 레이블 ──
\renewcommand{\figurename}{그림}
\renewcommand{\tablename}{표}

\title{사물인터넷과 빅데이터}
\subtitle{3주차: 프로젝트 기획}
\author{청주교육대학교 컴퓨터교육과}
\date{}

\begin{document}

% ── 슬라이드 1: 표지 ──
\begin{frame}[plain]
\maketitle
\end{frame}

% ── 슬라이드 2: 학습 목표 ──
\begin{frame}{오늘의 학습 목표}
\begin{enumerate}
  \item 학교 현장에서 센서로 측정 가능한 프로젝트 주제를 \textbf{선정}한다.
  \item 프로젝트 주제에 적합한 센서를 \textbf{선택}하고 이유를 설명한다.
  \item 데이터 수집 계획서를 작성하여 수집 절차를 \textbf{구체화}한다.
  \item ``측정 가능한 질문''과 ``측정 불가능한 질문''을 \textbf{구분}한다.
\end{enumerate}
\end{frame}

% ── 슬라이드 3: 프로젝트 4단계 ──
\begin{frame}{프로젝트 4단계 --- 지금 어디?}
\begin{center}
\includegraphics[width=0.85\textwidth]{그림3-1_프로젝트4단계흐름도.pdf}
\end{center}
\vfill
\begin{alertblock}{}
오늘 만드는 계획서 = 나머지 10주의 설계도
\end{alertblock}
\end{frame}

% ── 슬라이드 4: 좋은 주제 조건 ──
\begin{frame}{좋은 프로젝트 주제의 조건}
\begin{center}
\includegraphics[width=0.65\textwidth]{그림3-2_좋은주제4조건.pdf}
\end{center}
\begin{columns}[T]
\column{0.5\textwidth}
\textbf{필수 조건:}
\begin{itemize}
  \item 센서로 측정 가능
  \item 학교에서 수집 가능
\end{itemize}
\column{0.5\textwidth}
\textbf{권장 조건:}
\begin{itemize}
  \item 비교{\textperiodcentered}분석 가능
  \item 교육적 의미
\end{itemize}
\end{columns}
\end{frame}

% ── 슬라이드 5: 조건 1 ──
\begin{frame}{조건 1 --- 센서로 측정 가능한가?}
\textbf{센서 = 숫자를 만들어내는 장치}

\vspace{4mm}
\begin{columns}[T]
\column{0.48\textwidth}
\textbf{측정 가능}
\begin{itemize}
  \item 온도 (25.3℃)
  \item 소음 크기 (72)
  \item 밝기 (350 lux)
  \item \COtwo{} 농도 (800 ppm)
\end{itemize}

\column{0.48\textwidth}
\textbf{측정 불가능}
\begin{itemize}
  \item 학생의 기분
  \item 수업의 재미
  \item 급식의 맛
  \item 교실의 분위기
\end{itemize}
\end{columns}

\vfill
\begin{block}{핵심}
``숫자로 표현할 수 있는가?''
\end{block}
\end{frame}

% ── 슬라이드 6: 조건 2~4 ──
\begin{frame}{나머지 3가지 조건}
\textbf{조건 2: 학교에서 수집 가능한가?}\\
교실, 복도, 급식실, 도서관, 운동장 (야외 하천, 산속은 X)

\vspace{4mm}
\textbf{조건 3: 비교{\textperiodcentered}분석이 가능한가?}\\
시간 비교 (1교시 vs 5교시), 장소 비교 (1층 vs 3층), 조건 비교 (창문 열림 vs 닫힘)

\vspace{4mm}
\textbf{조건 4: 교육적 의미가 있는가?}\\
초등학교 현장에 적용 가능한 주제 $\rightarrow$ 미래 수업 설계에 도움
\end{frame}

% ── 슬라이드 7: 주제 카드 1/2 ──
\begin{frame}{프로젝트 주제 예시 카드 (1/2)}
\begin{center}
\small
\begin{tabular}{clll}
\toprule
& \textbf{주제} & \textbf{센서} & \textbf{난이도} \\
\midrule
1 & 교실 쾌적도 측정 & SHT40 + SGP30 & 중 \\
2 & 복도 소음 패턴 & 내장 마이크 & 하 \\
3 & 급식실 대기 환경 & VL53L0X + SHT40 & 중 \\
4 & 운동장 활동 환경 & SHT40 + TSL2561 & 중 \\
\bottomrule
\end{tabular}
\end{center}
\vfill
\begin{center}
\includegraphics[width=0.9\textwidth]{그림3-3_주제카드8종.pdf}
\end{center}
\end{frame}

% ── 슬라이드 8: 주제 카드 2/2 ──
\begin{frame}{프로젝트 주제 예시 카드 (2/2)}
\begin{center}
\small
\begin{tabular}{clll}
\toprule
& \textbf{주제} & \textbf{센서} & \textbf{난이도} \\
\midrule
5 & 도서관 조용함 & 내장 마이크 + 빛 센서 & 하 \\
6 & 화장실 사용 패턴 & VL53L0X & 중 \\
7 & 교실 환기 알림 & SGP30 + SHT40 & 상 \\
8 & 자연광 변화 관찰 & TSL2561 & 하 \\
\bottomrule
\end{tabular}
\end{center}

\vfill
\begin{block}{}
이 중에서 골라도, 변형해도, 새로 만들어도 OK!
\end{block}
\end{frame}

% ── 슬라이드 9: 측정 가능 vs 불가능 ──
\begin{frame}{``측정 불가능''을 ``측정 가능''으로 바꾸기}
\begin{center}
\includegraphics[width=0.8\textwidth]{그림3-4_측정가능vs불가능.pdf}
\end{center}

\vspace{2mm}
\small
\begin{center}
\begin{tabular}{lcl}
\toprule
\textbf{측정 불가능} & & \textbf{측정 가능} \\
\midrule
``덥다'' & $\rightarrow$ & ``온도가 28℃ 이상'' \\
``시끄럽다'' & $\rightarrow$ & ``소음이 70 이상'' \\
``쾌적하다'' & $\rightarrow$ & ``22{\tld}26℃, 40{\tld}60\% 시간 비율'' \\
``맛있다'' & $\times$ & (센서 변환 불가) \\
\bottomrule
\end{tabular}
\end{center}
\end{frame}

% ── 슬라이드 10: 핵심 메시지 ──
\begin{frame}{이번 학기에서 가장 중요한 메시지}
\begin{center}
\includegraphics[width=0.7\textwidth]{그림3-12_센서데이터기준판단흐름도.pdf}
\end{center}

\vfill
\begin{alertblock}{핵심}
\begin{itemize}
  \item \textbf{센서는 현재 상태만 측정한다.} 미래를 예측하지 않는다.
  \item \textbf{AI는 판단하지 않는다.} 우리가 만든 기준을 계산할 뿐이다.
\end{itemize}
\end{alertblock}
\end{frame}

% ── 슬라이드 11: 내장 센서 ──
\begin{frame}{마이크로비트 V2 내장 센서 (5종)}
\begin{columns}[T]
\column{0.45\textwidth}
\includegraphics[width=\textwidth]{그림3-5_마이크로비트V2센서배치도.pdf}

\column{0.52\textwidth}
\small
\begin{tabular}{lll}
\toprule
\textbf{센서} & \textbf{출력} & \textbf{특징} \\
\midrule
온도 & --5{\tld}50℃ & CPU +2{\tld}3℃ \\
빛 & 0{\tld}255 & 상대값 \\
마이크 & 0{\tld}255 & V2 전용 \\
가속도 & $\pm$2g & 3축 \\
나침반 & 0{\tld}360° & 자기장 \\
\bottomrule
\end{tabular}

\vspace{3mm}
{\footnotesize\color{coorange} 내장 온도: CPU 발열로 +2{\tld}3℃ 높음}
\end{columns}
\end{frame}

% ── 슬라이드 12: Grove 확장 센서 ──
\begin{frame}{Grove 확장 센서 --- I2C 연결}
\begin{columns}[T]
\column{0.45\textwidth}
\includegraphics[width=\textwidth]{그림3-6_GroveI2C연결도_3주차용.pdf}

\column{0.52\textwidth}
\small
\begin{tabular}{llll}
\toprule
\textbf{센서} & \textbf{모듈} & \textbf{측정} & \textbf{I2C} \\
\midrule
온습도 & SHT40 & $\pm$0.2℃ & 0x44 \\
공기질 & SGP30 & \COtwo{} & 0x58 \\
거리 & VL53L0X & mm & 0x29 \\
조도 & TSL2561 & lux & 0x39 \\
기압 & BMP280 & hPa & 0x76 \\
\bottomrule
\end{tabular}

\vspace{3mm}
{\footnotesize I2C = 센서마다 고유 ``주소''로 통신}
\end{columns}
\end{frame}

% ── 슬라이드 13: 내장 vs 확장 비교 ──
\begin{frame}{내장 센서 vs 확장 센서 비교}
\begin{center}
\includegraphics[width=0.85\textwidth]{그림3-13_내장vs확장센서비교.pdf}
\end{center}

\vfill
\begin{block}{}
처음이라면 내장 센서부터 시작하는 것도 좋아요!
\end{block}
\end{frame}

% ── 슬라이드 14: 도구 3종 ──
\begin{frame}{센서 선택 도구 3종}
\begin{center}
\includegraphics[width=0.85\textwidth]{그림3-14_도구연계흐름도.pdf}
\end{center}

\vspace{3mm}
\begin{center}
\includegraphics[width=0.9\textwidth]{그림3-8_도구3종화면.png}
\end{center}
\end{frame}

% ── 슬라이드 15: 실습 시뮬레이터 ──
\begin{frame}{[실습] 센서 시뮬레이터 v9 \hfill {\normalsize 5분}}
\begin{columns}[T]
\column{0.55\textwidth}
\includegraphics[width=\textwidth]{그림3-7_시뮬레이터v9화면.png}

\column{0.42\textwidth}
\textbf{실습 순서:}
\begin{enumerate}
  \item 시뮬레이터 열기
  \item 센서 드래그 $\rightarrow$ 포트 드롭
  \item 측정값 범위 확인
  \item 조합 메모
\end{enumerate}
\end{columns}
\end{frame}

% ── 슬라이드 16: 실습 카탈로그+검색 ──
\begin{frame}{[실습] 카탈로그 + 검색 도구 \hfill {\normalsize 7분}}
\textbf{카탈로그 (3분):} Excel 필터로 측정 항목 $\rightarrow$ 난이도 $\rightarrow$ 조합 확인

\textbf{검색 도구 (4분):} 키워드 입력 $\rightarrow$ 센서 추천

\vspace{3mm}
\begin{center}
\small
\begin{tabular}{ll}
\toprule
\textbf{키워드} & \textbf{추천 센서} \\
\midrule
``교실 환경'' & SHT40, SGP30 \\
``소음'' & 내장 마이크 \\
``밝기 정밀'' & TSL2561 \\
``사람 감지'' & VL53L0X \\
\bottomrule
\end{tabular}
\end{center}
\end{frame}

% ── 슬라이드 17: 의사결정 가이드 ──
\begin{frame}{센서 선택 의사결정 가이드}
\begin{center}
\includegraphics[width=\textwidth]{그림3-9_센서선택플로차트.pdf}
\end{center}
\end{frame}

% ── 슬라이드 18: 계획서란? ──
\begin{frame}{데이터 수집 계획서 = 프로젝트의 설계도}
\begin{columns}[T]
\column{0.45\textwidth}
\includegraphics[width=\textwidth]{그림3-10_수집계획서양식.pdf}

\column{0.52\textwidth}
\textbf{10개 필수 항목:}
\begin{enumerate}
  \item 프로젝트 제목
  \item 연구 질문
  \item 측정 항목
  \item 사용 센서
  \item 수집 장소
  \item 수집 간격
  \item 수집 기간
  \item 예상 데이터 양
  \item 예상 결과
  \item 주의 사항
\end{enumerate}
\end{columns}
\end{frame}

% ── 슬라이드 19: 연구 질문 ──
\begin{frame}{좋은 연구 질문 만들기}
\small
\begin{center}
\begin{tabular}{p{4cm}p{7cm}}
\toprule
\textbf{나쁜 예} & \textbf{좋은 예} \\
\midrule
``교실이 쾌적한가?'' & ``1교시와 5교시 중 쾌적 범위(22{\tld}26℃)에 있는 시간은 언제가 더 긴가?'' \\
``복도가 시끄러운가?'' & ``쉬는 시간 복도 소음은 수업 시간 대비 몇 배인가?'' \\
``밝기가 변하는가?'' & ``교실 창가 밝기가 가장 높은 시간대는 언제인가?'' \\
\bottomrule
\end{tabular}
\end{center}

\vfill
\begin{block}{포인트}
``{\tld}할까?'' 형태 + 센서 측정값 + 구체적 비교
\end{block}
\end{frame}

% ── 슬라이드 20: 수집 간격 ──
\begin{frame}{수집 간격 --- 얼마나 자주 기록할까?}
\begin{center}
\includegraphics[width=0.75\textwidth]{그림3-11_수집간격별데이터양.pdf}
\end{center}

\vspace{2mm}
\small
\begin{center}
\begin{tabular}{lccc}
\toprule
\textbf{대상} & \textbf{속도} & \textbf{권장} & \textbf{6시간} \\
\midrule
온도{\textperiodcentered}습도 & 느림 & 5분 & 72개 \\
\COtwo{} & 느림 & 5분 & 72개 \\
소음 & 빠름 & 10초{\tld}1분 & 360{\tld}2,160 \\
거리 & 매우빠름 & 1초 & 21,600 \\
\bottomrule
\end{tabular}
\end{center}
\end{frame}

% ── 슬라이드 21: 데이터 양 계산 ──
\begin{frame}{예상 데이터 양 계산}
\begin{block}{공식}
예상 데이터 개수 = (하루 수집 시간 $\div$ 수집 간격) $\times$ 수집 일수
\end{block}

\vspace{4mm}
\begin{center}
\includegraphics[width=0.7\textwidth]{그림3-15_저장용량시뮬레이션.pdf}
\end{center}

\vfill
{\footnotesize\color{coorange} 저장 용량 한계: 약 10,000개 (64KB). 1초 $\times$ 6시간 = 21,600개 $\rightarrow$ 초과!}
\end{frame}

% ── 슬라이드 22: 작성 예시 ──
\begin{frame}{작성 예시: ``교실 쾌적도 측정''}
\small
\begin{center}
\begin{tabular}{lp{8cm}}
\toprule
\textbf{항목} & \textbf{내용} \\
\midrule
제목 & 교실 쾌적도 시간대별 변화 측정 \\
연구 질문 & 1교시와 5교시 중 쾌적 범위에 있는 시간은? \\
측정/센서 & 온도, 습도, \COtwo{} / SHT40 + SGP30 \\
장소 & 3-1 교실 창가 책상 위 (70cm) \\
간격/기간 & 5분 / 5일(월{\tld}금) 09:00{\tld}15:00 \\
데이터 양 & 72개/일 $\times$ 5일 = 360개 \\
예상 결과 & 오후가 온도{\textperiodcentered}\COtwo{} 높을 것 \\
\bottomrule
\end{tabular}
\end{center}
\end{frame}

% ── 슬라이드 23: 실습 계획서 ──
\begin{frame}{[실습] 데이터 수집 계획서 작성 \hfill {\normalsize 12분}}
\textbf{순서:}
\begin{enumerate}
  \item 주제 카드 8종 또는 나만의 주제에서 선택
  \item 도구 3종으로 센서 결정
  \item 양식 10개 항목 채우기
\end{enumerate}

\vspace{4mm}
\textbf{자주 묻는 질문:}

\begin{description}
  \item[Q. 두 개 중 못 고르겠어요] 측정 항목이 겹치면 합칠 수 있어요
  \item[Q. 장소를 아직 못 정했어요] 5주차에 정해도 됩니다
  \item[Q. 예상 결과를 모르겠어요] 틀려도 OK!
\end{description}
\end{frame}

% ── 슬라이드 24: 주의 사항 ──
\begin{frame}{현장 수집에서 자주 생기는 문제}
\begin{center}
\begin{tabular}{ll}
\toprule
\textbf{문제} & \textbf{대비책} \\
\midrule
학생이 기기를 건드림 & 투명 케이스, 책꽂이 뒤 \\
배터리 방전 & 5분 간격, 매일 교체 \\
데이터 유실 & 매일 USB로 CSV 다운로드 \\
센서 오작동 & 수집 전 5분 테스트 \\
비{\textperiodcentered}먼지 (운동장) & 지퍼백 + 센서 구멍 노출 \\
\bottomrule
\end{tabular}
\end{center}

\vfill
\begin{block}{}
미리 생각해두면 5{\tld}8주차가 편해요!
\end{block}
\end{frame}

% ── 슬라이드 25: 발표 가이드 ──
\begin{frame}{계획서 간단 발표 (2{\tld}3명)}
\textbf{발표 시 반드시 포함할 3가지:}

\begin{enumerate}
  \item \textbf{연구 질문} --- ``나는 {\tld}를 알고 싶다''
  \item \textbf{센서 선택} --- ``{\tld}센서를 사용할 것이다. 왜냐하면 {\tld}''
  \item \textbf{예상 결과} --- ``{\tld}일 것이라고 예상한다''
\end{enumerate}

\vspace{5mm}
\textbf{동료 피드백 포인트:}
\begin{itemize}
  \item 센서로 측정 가능한 질문인가?
  \item 연구 질문에 적합한 센서인가?
  \item 학교에서 실제로 수집할 수 있는가?
\end{itemize}
\end{frame}

% ── 슬라이드 26: 과제 ──
\begin{frame}{이번 주 과제}
\begin{center}
\begin{tabular}{ll}
\toprule
\textbf{항목} & \textbf{내용} \\
\midrule
과제명 & 데이터 수집 계획서 \\
제출물 & 수집 계획서 (양식) + 개발일지 (PDF) \\
기한 & 일요일 자정 (LMS) \\
배점 & 5\% \\
\bottomrule
\end{tabular}
\end{center}

\vspace{4mm}
\textbf{개발일지 포함 내용:}
\begin{itemize}
  \item 오늘 한 것: 주제 선정 과정, 센서 선택 이유
  \item 고민한 것: 후보 주제 비교, 센서 조합
  \item 다음에 할 것: 4주차 수집기 제작 계획
  \item 느낀 점: 기획 경험 소감
\end{itemize}
\end{frame}

% ── 슬라이드 27: 등급 기준 ──
\begin{frame}{등급 기준}
\begin{center}
\begin{tabular}{ccl}
\toprule
\textbf{등급} & \textbf{점수} & \textbf{기준} \\
\midrule
A & 100\% & 질문 명확 + 센서 타당 + 계획 구체 + 일지 충실 + 통찰 \\
B & 85\% & 대부분 항목 작성 + 일지 충실 \\
C & 70\% & 일부 미흡 또는 일지 기본 수준 \\
D & 55\% & 다수 항목 미흡 \\
F & 0\% & 미제출 \\
\bottomrule
\end{tabular}
\end{center}

\vfill
\begin{block}{}
핵심: 연구 질문이 명확하고, 센서 선택에 이유가 있으면 높은 점수!
\end{block}
\end{frame}

% ── 슬라이드 28: 핵심 요약 ──
\begin{frame}{오늘의 핵심 요약}
\begin{enumerate}
  \item 좋은 주제 = \textbf{측정 가능} + \textbf{학교 수집} + 비교 가능 + 교육적 의미
  \item ``측정 불가능''을 ``측정 가능''으로 바꾸는 것이 핵심
  \item 도구 3종: 검색 $\rightarrow$ 카탈로그 $\rightarrow$ 시뮬레이터
  \item 계획서 10개 항목을 꼼꼼히 작성
  \item[\textbf{5.}] \textbf{센서는 현재 상태만 측정한다. 판단과 예측은 사람의 몫이다.}
\end{enumerate}
\end{frame}

% ── 슬라이드 29: 다음 주 예고 ──
\begin{frame}{4주차 예고: 데이터수집기 제작}
오늘 만든 계획서 $\rightarrow$ 다음 주의 설계도!

\vspace{5mm}
\begin{center}
\begin{tabular}{ll}
\toprule
\textbf{항목} & \textbf{내용} \\
\midrule
주제 & 데이터수집기 제작 \\
핵심 활동 & MakeCode 블록 코딩으로 Data Logger 작성 \\
사용 도구 & MakeCode 온라인 에디터 + 마이크로비트 V2 \\
과제 & MakeCode 파일 (.hex) + 테스트 결과 + 개발일지 \\
\bottomrule
\end{tabular}
\end{center}
\end{frame}

\end{document}
